\documentclass[12pt]{article}

\usepackage[margin=1.25in]{geometry}
\usepackage{setspace}
\usepackage{hyperref}
\usepackage{natbib}
\usepackage{amsmath}
\usepackage{amssymb}
\usepackage{microtype}
\usepackage{booktabs}
\usepackage{xcolor}
\usepackage{parskip}

\hypersetup{
  colorlinks=true,
  linkcolor=blue!60!black,
  citecolor=blue!60!black,
  urlcolor=blue!60!black
}

\onehalfspacing

\title{\textbf{State Bills of Credit vs.\ Continental Dollars:\\
Revolutionary War Finance and the Fiscal Theory of Paper Money}}

\author{Thomas J.\ Sargent}
\date{March 2026}

\begin{document}

\maketitle

\begin{abstract}
\noindent
During the American War of Independence (1775--1783), two distinct classes of
paper money circulated simultaneously: bills of credit issued by the thirteen
individual state governments and Continental dollars issued by the Continental
Congress.  Drawing on the work of Ferguson, Grubb, and others, this essay
documents two empirical regularities noted by Rothbard and supported by the
broader scholarly literature: (1)~state-issued currencies depreciated
significantly less than Continental dollars, and (2)~holders of state
currencies therefore earned higher ex-post real returns than holders of
Continental dollars.  The root cause of both differences is the structural
inability of the Continental Congress to levy taxes directly, which destroyed
the fiscal backing that gives a bill of credit its value.  Hamilton's Assumption
of 1790 covered state-issued \emph{bonds}, not state \emph{paper money};
the paper-money question had already been resolved, state by state, through
wartime taxation and post-war redemption programs.  The episode represents a
clean natural experiment in what is now called the fiscal theory of the price
level.
\end{abstract}

\tableofcontents

\newpage

%% ============================================================
\section{Introduction}
%% ============================================================

One of the relatively neglected episodes in early American fiscal history is
the simultaneous circulation of two quite different classes of paper money
during the Revolutionary War.  Each of the thirteen state governments
issued its own ``bills of credit''---paper instruments that served as local
currency---while the Continental Congress issued Continental dollars to
finance common war expenditures.  Both were legal tender (in varying and
contested degrees), both depreciated in terms of specie, and both ultimately
wound up being retired or repudiated in the post-war financial settlement.  Yet
the two classes behaved very differently.

Murray Rothbard, in volume~4 of \textit{Conceived in Liberty}
\citep{Rothbard1979}, argues that state currencies fared substantially better
than Continentals.  This claim is well-supported by the quantitative literature
inaugurated by \citet{Ferguson1961} and extended by \citet{Grubb2008,Grubb2012,
Grubb2003}.  The present essay documents that evidence, explains the
institutional mechanisms driving the differential, traces the variation across
states, and clarifies the relationship between state paper money and Hamilton's
Assumption plan of 1790.

The episode is an early and vivid illustration of what modern economists call
the \emph{fiscal theory of the price level}: a paper instrument is worth
something if and only if the issuing government credibly commits to levy
future taxes sufficient to redeem it.  The Continental Congress's
constitutional inability to tax directly is the proximate cause of the
Continental dollar's collapse to near-zero value; the states' retention of
direct taxing authority is the proximate cause of their currencies'
relative stability.

%% ============================================================
\section{The Institutional Setting}
%% ============================================================

\subsection{Two Issuers, Two Instruments}

Congress began issuing Continental dollars in June 1775.  Total face-value
emissions reached approximately \$241~million by 1779, when Congress largely
ceased new issues.  By 1780--81 the Continental had depreciated to roughly
100~paper dollars per specie dollar, and Congress formally acknowledged
the collapse with its ``40-to-1'' devaluation resolution of March~1780.  The
phrase ``not worth a Continental'' entered common speech as a permanent record
of the episode \citep{Ferguson1961, Timberlake1993}.

At the same time, each state issued its own bills of credit in varying amounts.
These circulated alongside Continentals, were legally distinct obligations
of the individual state governments, and were redeemable (in theory) through
state tax receipts.  The aggregate volume of state emissions was smaller than
the Congressional total, and the depreciation experience varied considerably
across states but was uniformly less severe than the Continental's collapse.

\subsection{What Is a Bill of Credit?}

A ``bill of credit'' is essentially a government promise: the issuer pledges
to accept the instrument in payment of future taxes.  Its market value
reflects the credibility of that pledge---the probability that the government
will in fact levy sufficient taxes and honor the redemption
\citep{Smith1985, Grubb2016}.  This is precisely the mechanism that modern
fiscal theories of the price level formalize: the real value of outstanding
paper equals the present discounted value of future primary surpluses
available for its redemption \citep{Perkins1994}.

%% ============================================================
\section{Claim~1: State Currencies Depreciated Less}
%% ============================================================

\subsection{The Quantitative Evidence}

The most careful quantitative treatment is that of \citet{Grubb2008,Grubb2012},
who assembles exchange-rate data against specie for both Continental dollars
and individual state currencies, and constructs estimates of real returns
to holders.  The central finding:

\begin{itemize}
  \item Continental dollars fell to approximately \textbf{100:1} against specie
    by 1780--81 before Congress stopped issuing them.
  \item State currencies depreciated by factors ranging from roughly \textbf{3:1
    to 20:1} in most cases---severe, but far less catastrophic than the
    Continental.
  \item The best-performing state currencies (Massachusetts, Connecticut) held
    their value at roughly \textbf{3--5:1} against specie at the nadir.
  \item The worst-performing state currencies (North Carolina, Georgia, South
    Carolina) reached 40:1 or worse, largely owing to British military
    occupation that paralyzed state tax administration.
\end{itemize}

This pattern is also documented in \citet{Ferguson1961} and
\citet{Sumner1968}.

\subsection{Why Did State Currencies Fare Better?}

Five institutional factors drove the differential:

\paragraph{1.  Tax-backing credibility.}
States possessed direct, constitutionally unquestioned taxing authority over
their own residents.  The Continental Congress, under the Articles of
Confederation, could only \emph{requisition} funds from states---requests
that were routinely delayed or ignored.  A bill of credit is only as good as
the taxing power behind it; the states had that power, the Congress did not
\citep{Ferguson1961, Edling2014}.

\paragraph{2.  Volume discipline.}
Individual state emissions were smaller in absolute terms than the
Congressional total.  For a given tax base, a smaller emission implies a
smaller required future surplus and therefore a more credible redemption
pledge \citep{Grubb2003}.

\paragraph{3.  Geographic compactness.}
Tax assessment and collection were more tractable within a single state than
across a continent.  Enforcement capacity reinforced the credibility of the
redemption pledge.

\paragraph{4.  Competitive discipline.}
States competed for residents, commerce, and reputational standing.  A state
that destroyed its currency harmed itself directly and visibly in ways that
the Congress, lacking a direct constituency, did not.

\paragraph{5.  Free-rider externality.}
The Continental Congress suffered from a classic collective-action problem:
each state benefited from common defense financed by Continental inflation,
while the inflation tax fell across all holders of Continentals.  Individual
states internalizing the costs of their own emissions did not face this
externality in the same form \citep{MatsonOnuf1990}.

%% ============================================================
\section{Claim~2: Holders of State Currencies Earned Higher Real Returns}
%% ============================================================

\subsection{The Return Comparison}

The terminal redemption experience of the two instruments was starkly
different.

\textbf{Continental dollars.}  After the collapse, Hamilton's First Report on
Public Credit (January 1790) proposed---and Congress enacted in the Funding
Act of 1790---redeeming residual Continentals at the rate of
\textbf{100~Continental dollars = 1~new federal dollar} in long-term federal
bonds.  This made explicit the roughly 100:1 specie depreciation that markets
had already registered.  Holders who had retained Continentals from early
emissions received approximately 1~percent of original face value in real
terms.

\textbf{State bills of credit.}  Most states had been actively retiring their
own currencies through earmarked taxation throughout and immediately after the
war.  The depreciation having been less severe---3:1 to 20:1 vs.\ 100:1---the
implied real loss to holders was much smaller.  In New England, where tax
programs were most faithfully executed, some holders received redemptions close
to original specie value \citep{Grubb2012}.

The Rothbard claim---that ex-post real returns to state currency holders
exceeded those to Continental holders---therefore holds by construction: it
follows directly from the lesser depreciation during the war, since the
post-war redemption rates in both cases were broadly in line with wartime
market prices.

\subsection{Caveats}

Two qualifications are important.  First, the comparison is of a
\emph{weighted average} across state currencies: holders in Georgia or South
Carolina did nearly as poorly as Continental holders.  Second, the exact
return depends on the timing of holding: a speculator who bought depreciated
Continentals at 80:1 and redeemed them at the official 100:1 rate suffered a
further loss of 20~percent, but one who bought at 95:1 and sold at 100:1
suffered only marginally.  Grubb's detailed calculations account for these
timing effects.

%% ============================================================
\section{Hamilton's Assumption and Its Scope}
%% ============================================================

A natural question is whether Hamilton's Assumption plan of 1790 covered
state-issued paper money.  The answer is \textbf{no}.

\subsection{What Assumption Covered}

The Funding Act of 1790 assumed approximately \$21.5~million in state
\emph{war debts}---the interest-bearing certificates, soldiers' pay warrants,
and loan-office notes that states had issued as explicit debt instruments to
finance war expenditures.  These were recognized obligations to identifiable
creditors \citep{Ferguson1961, Edling2014, Perkins1994}.

\subsection{What Assumption Did Not Cover}

State \textit{bills of credit} used as circulating medium were not part of the
Assumption for two reasons:
\begin{enumerate}
  \item Most had already been retired through taxation or had become worthless
    in the market by 1790.  There was no meaningful stock of state paper
    money outstanding at the time of the Funding Act.
  \item They were not legal obligations of the federal government in any
    constitutional sense.  The state paper-money question was handled
    entirely within each state's own fiscal and legal framework.
\end{enumerate}

Continentals were treated as a federal obligation (redeemed at the 100:1
rate), but this was a special case arising from their status as instruments
of the Congress itself.

\subsection{The Political Economy of Assumption}

The assumption controversy---the famous Jefferson-Madison-Hamilton dinner
bargain that located the permanent capital on the Potomac---was driven
entirely by the distribution of state \emph{bond} debts.  Virginia had
substantially retired its war bonds through taxation; Virginia therefore
opposed Assumption as a transfer to less fiscally responsible states
(especially South Carolina and Massachusetts, which had large residual
balances).  South Carolina and Georgia, whose fiscal administrations had
been devastated by British occupation, had the largest outstanding state
bond debts and therefore gained most from Assumption \citep{Ferguson1961,
MatsonOnuf1990, Edling2014}.

%% ============================================================
\section{Cross-State Variation and Fiscal Policy}
%% ============================================================

The differential performance of state currencies depended directly on specific
fiscal choices made by each state government.  Table~\ref{tab:states} summarises
the broad patterns.

\begin{table}[htbp]
\centering
\caption{State Currency Performance and Fiscal Policies, 1775--1783}
\label{tab:states}
\begin{tabular}{llp{7cm}}
\toprule
\textbf{State} & \textbf{Currency Performance} & \textbf{Key Fiscal Factors} \\
\midrule
Massachusetts   & Relatively good (3--5:1)  & Early aggressive taxation; sinking-fund
                                              earmarks; limited total emission \\[4pt]
Connecticut     & Relatively good (3--5:1)  & Annual tax levies earmarked for currency
                                              redemption; disciplined emission control \\[4pt]
Pennsylvania    & Moderate (5--10:1)        & Larger emissions but significant tax effort;
                                              loan-office structure provided partial backing \\[4pt]
New York        & Moderate (6--12:1)        & Interrupted by British occupation of
                                              Manhattan; recovered as tax collections
                                              resumed \\[4pt]
Virginia        & Significant (15--30:1)    & Large tobacco-economy; difficult specie
                                              valuation; ultimately redeemed at
                                              unfavorable rates \\[4pt]
North Carolina  & Severe (30--50:1)         & Large emissions relative to small tax base;
                                              chaotic wartime administration \\[4pt]
South Carolina  & Severe to catastrophic    & British occupation 1780--82 devastated tax
                                              collection; currency near-worthless; large
                                              bond debt assumed by Hamilton \\[4pt]
Georgia         & Severe to catastrophic    & Same factors as South Carolina; smallest
                                              and most vulnerable state \\
\bottomrule
\end{tabular}
\end{table}

\subsection{Tax-Earmarking and Redemption Plans}

The mechanism linking fiscal policy to currency value is clearest in the
New England cases.  Connecticut, for example, legislated specific annual
tax levies at the outset of each currency emission, credibly pre-committing
the state's tax revenues to future redemption.  This pre-commitment is
precisely what the fiscal theory of the price level requires to stabilize
currency value: agents holding the bills knew that future taxes would
be levied, creating demand for the bills as a means of tax payment
\citep{Grubb2016, Smith1985}.

Pennsylvania took an intermediate approach, relying partly on a loan-office
system under which bills were lent to private borrowers collateralized by
land, and partly on direct tax levies.  This blended structure provided
partial but imperfect backing, consistent with intermediate depreciation
\citep{Perkins1994}.

The southern states combined large emissions with weak tax-collection
capacity even before British occupation.  Virginia's currency depreciated
more than might be expected from its fiscal fundamentals partly because of
a chronic shortfall between legislated taxes and actual collections.

\subsection{Debt Rescheduling}

Several states converted near-term currency liabilities into longer-dated
interest-bearing bonds---essentially conducting a debt exchange from
floating paper to funded debt.  This operation did not eliminate the fiscal
burden but it extended the maturity profile, reducing the immediate demands
on tax revenue and helping to stabilize currency value in the short run.
The post-war funded debts created by this process were exactly the instruments
that Hamilton's Assumption absorbed \citep{Ferguson1961}.

%% ============================================================
\section{The Fiscal Theory Interpretation}
%% ============================================================

\citet{Grubb2008, Grubb2012, Grubb2003} explicitly frames the Continental
dollar and state currency episodes in terms of what modern economists call
the \emph{fiscal theory of the price level}.  The real value of a stock of
outstanding government paper is determined by the present discounted value
of future primary fiscal surpluses committed to its redemption.

Under this interpretation the story of the Revolutionary War currencies
reduces to a comparison of fiscal backing:

\begin{align*}
  V_t &= \mathbb{E}_t \left[ \sum_{s=0}^{\infty} \beta^s \, S_{t+s} \right]
\end{align*}

where $V_t$ is the real value of outstanding paper at time~$t$, $S_{t+s}$
is the real primary surplus at $t+s$ earmarked for redemption, and $\beta$
is the appropriate discount factor.  For the Continental Congress,
$S_{t+s} \approx 0$ because the Congress had no independent taxing power;
the expected stream of surpluses was negligible, driving $V_t$ toward zero.
For states that credibly legislated tax plans, $S_{t+s}$ was positive and
observable, keeping $V_t$ well above zero.

This framing also clarifies why \emph{announcements} mattered so much.
When a state legislature voted a specific tax schedule tied to currency
redemption, it provided hard information about $\{S_{t+s}\}$, immediately
raising $V_t$.  When the Continental Congress printed more dollars without
any credible tax backing, it diluted $V_t$ across a larger stock without
increasing the backing, driving down the price level.

%% ============================================================
\section{Conclusion}
%% ============================================================

The Revolutionary War paper-money episode supports both of the claims
noted by Rothbard and examined here.  State-issued bills of credit
depreciated substantially less than Continental dollars, and holders of
state bills therefore earned higher ex-post real returns.  The fundamental
cause of the difference is fiscal: states possessed direct, credible taxing
authority; the Continental Congress did not.  States that made specific,
binding tax commitments to currency redemption (New England) performed best;
states whose tax administration collapsed under British occupation (South
Carolina, Georgia) performed worst.

Hamilton's Assumption of 1790 absorbed state \emph{bonds}, not state
\emph{paper money}.  The paper-money story had been resolved, state by
state, through wartime and post-war taxation before the Funding Act was
written.  The Assumption controversy was about the distribution of funded
bond debt---a politically charged transfer whose geography closely tracked
the severity of British military occupation.

The episode remains one of the clearest historical illustrations of the
fiscal theory of the price level: government paper is worth what the
issuing government's future taxes will redeem it for, no more and no less.

%% ============================================================
\bibliographystyle{plainnat}
\bibliography{revolutionary_war_state_currencies}
%% ============================================================

\end{document}
